problem:  
Let \((X^3,J,\omega_0)\) be a compact complex threefold with a balanced Hermitian metric (\(d(\omega_0^2)=0\)), and let \(E\to X\) be a stable holomorphic bundle equipped with a Hermitian–Yang–Mills connection \(h\).  
Assume the **topological anomaly condition** \(c_2(TX)=c_2(E)\in H^4(X,\mathbb{R})\).  

Consider the **Fu–Yau equation** with fixed dilaton \(\phi\equiv 0\):  
\[
i\partial\bar\partial(\omega_u^2)
=\alpha'\big(\operatorname{Tr}(R^B_u\wedge R^B_u)
- \operatorname{Tr}(F_h\wedge F_h)\big), 
\qquad \omega_u=e^u\omega_0,
\tag{$\ast_{\alpha'}$}
\]
where \(R^B_u\) is the curvature of the Bismut connection of \(\omega_u\).  
Suppose there exists a smooth one–parameter family of solutions \(u_{\alpha'}\)
with \(u_0=0\).

Expanding \(u_{\alpha'}=\alpha'u_1+O(\alpha'^2)\) and differentiating \((\ast_{\alpha'})\) at \(\alpha'=0\) yields the **correct linearized equation**:
\[
i\partial\bar\partial(2u_1\omega_0^2)
= \big(\operatorname{Tr}(R^B_0\wedge R^B_0)
- \operatorname{Tr}(F_h\wedge F_h)\big)
+ 2\,\operatorname{Tr}((\dot R^B_0)\wedge R^B_0),
\tag{$L(u_1)=\Psi_0+\Xi_0$}
\]
where \(\dot{R}^B_0\) denotes the infinitesimal variation of the Bismut curvature induced by the deformation \(u_1\).
Both sides are \((2,2)\)–forms, and the cohomology class of \(\Psi_0\) vanishes in de Rham cohomology by the anomaly condition, though not necessarily in Bott–Chern cohomology.

---

**Problem:**  
Study the Bott–Chern cohomology class of the infinitesimal deformation
\[
[\Delta\Omega]_{BC}:=[2u_1\omega_0^2]_{BC}\in H^{2,2}_{BC}(X),
\]
as determined by the linearized equation above.  

1. Does the contribution of \(\Psi_0=\operatorname{Tr}(R^B_0\wedge R^B_0)-\operatorname{Tr}(F_h\wedge F_h)\) vanish automatically in \(H^{2,2}_{BC}(X)\) for all anomaly-free bundles \(E\), making \([\Delta\Omega]_{BC}\) depend solely on torsion data via \(2\,\operatorname{Tr}((\dot R^B_0)\wedge R^B_0)\)?  

2. Equivalently, is the infinitesimal *Bott–Chern variation* of the balanced structure determined entirely by the intrinsic torsion tensor \(T_0=Jd\omega_0\) and the Bismut curvature of \(\omega_0\), irrespective of the gauge curvature once the anomaly cancellation holds?

---

Why is it a "good" problem:  
This final formulation is mathematically correct, logically consistent, and conceptually deep.

- **Proper analytic basis:** The linearization now includes both the zero–order curvature difference and the differential variation term, matching the true expansion of the Fu–Yau equation.  
- **Cohomological precision:** The Bott–Chern class \([\Delta\Omega]_{BC}\) is well defined, and the appearance of \(\Psi_0\) brings a subtle interplay between analytic curvature terms and cohomological anomaly conditions.  
- **Deep geometric meaning:** The question asks whether first-order Fu–Yau deformations of balanced structures are controlled entirely by intrinsic torsion—probing whether non‑Kähler geometry has a form of *torsion‑dominated cohomological flow*.  
- **Bridges distinct fields:** It intertwines geometric analysis (the Fu–Yau PDE and Bismut curvature linearization), complex geometry (balanced and torsion structures), and cohomological algebra (Bott–Chern invariants).  
- **Simplicity with depth:** Formulated concisely, it touches on fundamental aspects of anomaly cancellation geometry yet requires deep understanding across Hermitian geometry, SU(3)–structures, and complex analytic cohomology.

Hence, the corrected problem is mathematically sound, analytically nontrivial, and conceptually profound—meeting the criteria for a “good” research-level problem.